\documentclass[a4paper,11pt]{article}
\usepackage[utf8]{inputenc}
\usepackage{geometry}
\geometry{margin=2cm, top=2.5cm, bottom=2.5cm}
\usepackage{graphicx}
\usepackage{fancyhdr}
\usepackage{titlesec}
\usepackage[hidelinks]{hyperref}
\usepackage{caption}
\usepackage{setspace}
\usepackage{listings}
\usepackage{xcolor}
\usepackage{float}
\usepackage{enumitem}
\usepackage{longtable}
\usepackage{booktabs}
\usepackage{minted}
\usepackage{tcolorbox}

\setstretch{1.3}
\pagestyle{fancy}
\fancyhf{}
\fancyhead[L]{\textit{HolisticX Infrastructure: Developer Handover & Guide}}
\fancyfoot[C]{\thepage}

\definecolor{sectioncolor}{rgb}{0.1,0.1,0.4}
\definecolor{boxback}{rgb}{0.97,0.97,0.97}

\titleformat{\section}
  {\normalfont\Large\bfseries\color{sectioncolor}}
  {\thesection}{1em}{}
\titleformat{\subsection}
  {\normalfont\large\bfseries\color{sectioncolor}}
  {\thesubsection}{1em}{}

\newtcolorbox{guidebox}[1]{
  colback=boxback,
  colframe=sectioncolor,
  fonttitle=\bfseries,
  title=#1
}

\title{\textbf{HolisticX Infrastructure: Developer Handover \& Guide}}
\author{Saifeddine Douidy / ToumAI}
\date{\today}

\begin{document}

\maketitle
\thispagestyle{empty}
\tableofcontents
\newpage

\section{Introduction \& Core Principles}
This document is the primary technical guide for the HolisticX cloud infrastructure. It details the Terraform code, the multi-environment deployment strategy, and the underlying AWS architecture.

\begin{itemize}
    \item \textbf{Infrastructure as Code (IaC):} All infrastructure is managed by Terraform. \textbf{Do not make manual changes in the AWS Console.}
    \item \textbf{Modularity:} The code is organized into reusable modules for clarity and maintainability.
    \item \textbf{Environment Isolation:} Each environment (dev, staging, prod) is a completely separate, isolated deployment managed by Terraform Workspaces, each with its own state file.
\end{itemize}

\section{Project Setup \& Prerequisites}
\begin{guidebox}{Action Required: Configure the Terraform Backend}
Before you can deploy anything, you must configure the remote backend. This is where Terraform will store its state file, which is critical for collaboration and state locking.

\begin{enumerate}
    \item \textbf{Create an S3 Bucket:} Manually create a private S3 bucket in your AWS account. This bucket must be created before you run Terraform. A good name would be \texttt{holisticx-terraform-state-<your-account-id>}.
    \item \textbf{Update Backend Configuration:} Open the file \texttt{live/my-project/backend.tf}. Update the \texttt{bucket} name to the one you just created.
    \begin{minted}{terraform}
# In live/my-project/backend.tf
terraform {
  backend "s3" {
    bucket = "holisticx-terraform-state-<your-account-id>" # <-- CHANGE THIS
    key    = "live/my-project/terraform.tfstate"
    region = "us-east-1"
  }
}
    \end{minted}
    \item \textbf{Initialize Terraform:} Run \texttt{terraform init}. This will configure your local setup to use the remote S3 bucket for state storage.
\end{enumerate}
\end{guidebox}

\section{Code Structure Explained}
\begin{description}
    \item[\texttt{modules/}] A library of reusable architectural components (VPC, RDS, etc.).
    \item[\texttt{live/my-project/}] The composition root. This is where you run all Terraform commands. It assembles the modules to build the complete infrastructure.
\end{description}

\section{Multi-Environment Strategy with Workspaces}
We use Terraform Workspaces to manage multiple environments from a single codebase.

\subsection{The Workspace-to-Environment Mapping}
The logic in \texttt{live/my-project/locals.tf} translates a flexible workspace name (e.g., \texttt{dev-saifeddine}) into a standardized environment tag (e.g., \texttt{Dev}) that is applied to all resources for cost tracking.

\begin{minted}{terraform}
locals {
  // This lookup finds the current workspace name in a map.
  // If not found, it defaults to "Dev". This is powerful because
  // any new developer workspace is automatically tagged correctly.
  environment_mapped = lookup(
    var.workspace_to_environment_tag_map, 
    terraform.workspace, 
    "Dev"
  )
}
\end{minted}

\section{Deployment Guide}
\subsection{Step 1: Create and Select Your Workspace}
\begin{minted}{bash}
# Run from live/my-project/
cd live/my-project

# Create a workspace for your feature or environment
terraform workspace new feature-billing

# Switch to it
terraform workspace select feature-billing
\end{minted}

\subsection{Step 2: Create Your Variable File}
Copy an existing \texttt{.tfvars} file (\texttt{cp env/develop.tfvars env/feature-billing.tfvars}) and modify it for your needs. \textbf{Crucially, you must provide a globally unique name for the \texttt{rds\_backup\_s3\_bucket\_name} variable.}

\subsection{Step 3: Cost Estimation (Best Practice)}
Before applying, always estimate the cost. This project is configured to use **Infracost**.
\begin{minted}{bash}
# Authenticate with Infracost Cloud (one-time setup)
infracost auth login

# Generate a cost breakdown based on your variable file
infracost breakdown --path . \
  --terraform-var-file="env/feature-billing.tfvars" \
  --usage-file="infracost-usage.yml"
\end{minted}
Review the output to understand the financial impact of your changes.

\subsection{Step 4: Plan and Apply}
\begin{minted}{bash}
# Always run a plan first to see the intended changes
terraform plan -var-file="env/feature-billing.tfvars"

# If the plan is correct, apply it
terraform apply -var-file="env/feature-billing.tfvars"
\end{minted}

\subsection{Step 5: Post-Deployment DNS}
For the first deployment, you must delegate DNS from your registrar (IONOS) to AWS.
\begin{minted}{bash}
# Get the AWS name servers for your new Route 53 Hosted Zone
terraform output name_servers
\end{minted}
Log in to IONOS and update the NS records for your domain to these four values.

\section{Future Enhancements: Terraform Cloud}
For enhanced collaboration, security, and governance, migrating to **Terraform Cloud** is the recommended next step.
\begin{itemize}
    \item \textbf{Centralized State Management:} It provides a more robust and secure alternative to the S3 backend.
    \item \textbf{Collaborative Runs:} It allows for team-based review and approval of Terraform runs.
    \item \textbf{Cost Estimation:} It can be integrated with Infracost to automatically show a cost estimate on every pull request.
    \item \textbf{Policy as Code:} Enforce security and compliance policies automatically.
\end{itemize}

\section{Architectural Reference}
For a high-level overview of the cloud architecture, including diagrams and security summaries, please refer to the \texttt{HolisticX\_Cloud\_Architecture\_Report.tex} document.

\end{document}